\documentclass[12pt,a4paper]{article}

% Required packages for math symbols and formatting
\usepackage{amsmath}
\usepackage{amssymb}
\usepackage{enumerate}
\usepackage[margin=1in]{geometry}

% Required package for Traditional Chinese support via pdflatex
\usepackage{CJKutf8}

\begin{document}
\begin{CJK*}{UTF8}{bsmi} % 'bsmi' is the standard Traditional Chinese font family in CJKutf8

\begin{center}
    \Large \textbf{Assignment}
\end{center}

\vspace{0.5cm}

\section*{Problem 1 (15\%)}
Give the double binary representation of the number $+17.3$.

\vspace{0.5cm}

\section*{Problem 2 (20\%)}
Let
\[ \beta=10, \quad p=4, \quad x=1.000, \quad y=-0.5555 \]
Consider
\begin{align*}
x_0 &= x \\
x_1 &= (x_0 \ominus y) \oplus y \quad \text{(*註一)} \\
&\vdots \\
x_n &= (x_{n-1} \ominus y) \oplus y
\end{align*}
Assume exactly rounded operations using the rounding up scheme. What are the values of the sequence $\{x_0, x_1, \dots\}$?

\vspace{0.5cm}

\section*{Problem 3 (25\%)}
Conduct LU factorization with pivoting for the following matrix
\[
A = \begin{bmatrix}
4 & -10 & -3 & 7 \\
6 & -5 & -4 & 7 \\
6 & 27 & 3 & -3 \\
12 & 6 & -6 & 6
\end{bmatrix}
\]

\begin{enumerate}[(a)]
    \item What are $M_3, P_3, M_2, P_2, M_1, P_1$ such that
    \[ M_3 P_3 M_2 P_2 M_1 P_1 = U \ ? \]
    
    \item What is
    \[ P_1^{-1} M_1^{-1} P_2^{-1} P_3^{-1} \ ? \]
    
    \item What is
    \[ P_1^{-1} M_1^{-1} P_2^{-1} M_2^{-1} P_3^{-1} M_3^{-1} \ ? \]
    
    \item What is
    \[ P_3 P_2 P_1 P_1^{-1} M_1^{-1} P_2^{-1} P_3^{-1} \ ? \]
    
    \item What is the $P$ such that
    \[ P A = L U \]
    Is it equal to $P_3 P_2 P_1$?
\end{enumerate}

\vspace{0.5cm}

\section*{Problem 4 (20\%)}
Consider a binary number in a floating-point system of $\beta=2$ and $p=6$:
\[ 1.11111 \]

\begin{enumerate}[(1)]
    \item Transform it to a decimal value with $\beta=10$ and $p=2$. That is,
    \[ x.x \]
    Then transform this decimal value to a binary representation with $p=7$. Is the resulting binary the same as (1)?
    
    \item Redo (1) but a decimal value with $p=3$.
\end{enumerate}

\noindent Note that we assume rounding even is used.

\vspace{0.5cm}

\section*{Problem 5 (20\%)}
In our lecture we derive a version of Cholesky factorization by
\[ L_{jj} L_{:j} = A_{:j} - \sum_{k=1}^{j-1} L_{jk} L_{:k} \]
By a similar setting please derive a version that obtains
\[ A = U^T U \quad \text{(*註二)} \]
You need to give the program.

\vspace{1cm}
\hrule
\vspace{0.5cm}

\noindent \textbf{Footnotes:} \\
*註一 : $\ominus$ 為內有減號的圈圈, 即 rounded subtraction. \\
*註二 : $U^T$ ($U'$) 為 transpose of upper triangular matrix $U$.

\end{CJK*}
\end{document}